%! Populations, Samples, Parameters, and Statistics — Unit 1, Lesson 1.2 — Lesson Plan
% PILOT SHAPE (2026-09-06): modeled on AP Statistics lesson 1.4 minus its AP Practice page.
% GRADUAL-RELEASE plan (warm-up / guided notes and practice / spoken debrief / close and START
% THE HOMEWORK). The guided notes carry the release ROW BY ROW in one Main Ideas / Notes table
% (2026-09-12 shape): I do / we do in rows 1-4 (problems 1-15) -> we do the Guided Practice row
% (problems 16-19). THERE IS NO INDEPENDENT PRACTICE SET IN THE NOTES — the homework is the individual
% practice, and the period ends with students starting it. Teacher-facing. Breakable boxes are
% `skillbox` with a \boxguard ahead; the two tabularx boxes are `fixedskillbox` (never splits).
% No group activity, no tier boxes, no exit ticket.
\documentclass[10pt]{article}
\usepackage{saar-article}
\usepackage{saar-boxes}
\usepackage{graphicx}   % -article does NOT load graphicx
\graphicspath{{images/}}
\newcommand{\TallMath}[1]{$\displaystyle #1\rule[-1.4em]{0pt}{3.2em}$}
\newcommand{\UnitNumberName}{Unit 1: Asking Questions with Data: Study Design \quad}
\newcommand{\LessonNumberName}{Lesson 1.2: Populations, Samples, Parameters, and Statistics}

\begin{document}
\begin{center}
    {\Large\bfseries \CourseName \\
    {\small\bfseries \UnitNumberName \LessonNumberName}}
\end{center}

\begin{tcolorbox}[colback=frost, colframe=cerulean, boxrule=0.9pt, arc=2mm,
    left=2mm, right=2mm, top=1.5mm, bottom=1.5mm]
    \textbf{Primary Objective:} Students will name the population and the sample in a study,
    label a given number as a parameter or a statistic by asking \emph{who it describes},
    name the real-world constraint that keeps the study from being a census, explain why two
    honest samples give two different statistics for one unchanging parameter, and judge
    what a sample does --- and does not --- allow a report to claim.
    \\[2pt]
    \textbf{Standards:} PS.DC.1d, PS.DC.1e \quad$\cdot$\quad PS.DC.1b carried forward in the
    conclusion items \quad$\cdot$\quad previews PS.DC.2a--c (Lessons 1.3--1.4)
    \\[2pt]
    \textbf{Lesson model:} \emph{Gradual release --- I do, we do, you do --- all of it inside
    the guided notes.} There is no group activity. The notes name the vocabulary as they teach
    it and deliver the instruction \textbf{row by row in one Main Ideas / Notes table} (four
    rows, problems 1--15: you read the definition, mark up the display, and work the first
    problem of each grid; the class works the rest); the Guided Practice row (problems 16--19)
    is worked together with students holding the pen; \textbf{the homework is the individual practice}, and students
    start it in class while you circulate. The whole-class debrief is \emph{spoken} --- there is
    no practice set at the end of the notes and no exit ticket.
    \\[2pt]
    \textbf{Note on notation:} $\mu$, $\bar{x}$, $p$, $\hat{p}$ are deliberately not on the
    page. If a student asks, say the symbols return in Units 4, 7, and 8 and move on --- do not
    introduce, assign, or test them today.
\end{tcolorbox}

\renewcommand{\arraystretch}{1.4}
\boxguard
\begin{skillbox}[Learning Targets \& Key Understandings]{goldbox}
\begin{tabularx}{\linewidth}{@{}X|X@{}}
\textbf{Learning targets (student language)} & \textbf{Key understandings (the ``why'')} \\[2pt]
\hline
\vspace{-6pt}
\begin{itemize}[leftmargin=1.1em, nosep, after=\vspace{2pt}]
  \item I can name the \textbf{population} and the \textbf{sample} in a study, and give the
        sample size. \hfill \textit{PS.DC.1d}
  \item I can tell a \textbf{parameter} from a \textbf{statistic} by asking who the number
        describes --- never by what it looks like. \hfill \textit{PS.DC.1d}
  \item I can name the \textbf{constraint} --- time, cost, access, destruction, change --- that
        keeps a study from being a census. \hfill \textit{PS.DC.1e}
  \item I can explain why two honest samples give two different statistics for one parameter,
        and say what a sample does and does not let a report claim.
        \hfill \textit{PS.DC.1d--e; 1b}
\end{itemize}
&
\vspace{-6pt}
\begin{itemize}[leftmargin=1.1em, nosep, after=\vspace{2pt}]
  \item \emph{The population is decided by the question, not by who is available.} Naming it
        is a stage-1 act; the sample is chosen out of it afterwards.
  \item \textbf{Target misconception:} that two samples disagreeing means somebody made a
        mistake --- or that the ``real'' answer is one of them. There is \textbf{one}
        parameter and it did not move; the three volunteers' $52\%$, $48\%$, and $58\%$ are
        three honest estimates of it. \emph{Sampling variability is what samples do.}
  \item \textbf{Second misconception:} that a percent is a statistic and a count is a
        parameter. The kind is decided by \emph{who the number describes}: $61\%$ drawn from
        the records of all $900$ students is a parameter.
  \item Constraints are not excuses. They are why sampling exists, and part of describing a
        study honestly (PS.DC.1e).
  \item A conclusion may describe only the population the sample was drawn from, and only
        the question that was asked --- change either and the same statistic estimates a
        different parameter, or nothing at all.
\end{itemize}
\end{tabularx}
\end{skillbox}

\boxguard[24]
\begin{skillbox}[{Vocabulary, Concepts, \& Theorems --- taught directly in the guided notes}]{greenbox}
\small
\renewcommand{\arraystretch}{1.35}
\begin{tabularx}{\linewidth}{@{} >{\bfseries}p{3.1cm} X @{}}
Population & The entire group the question is about --- every individual the conclusion should
       describe. \\
Sample & The part of the population data were actually collected from. \textbf{Sample size}
       ($n$) --- how many individuals are in it. \\
Census & A study that collects data from every individual in the population. \\
Parameter & A number that describes the population. Usually unknown --- and there is only one. \\
Statistic & A number computed from the sample. Known --- and it changes from sample to sample. \\
Constraint & A real-world limit on how a study can be run: time, cost, access, destruction, or
       a changing population. \\
Sampling variability & The change in a statistic from one sample to the next, even when the
       population has not changed. Not a mistake. \\
Estimate & Using a statistic as the best available stand-in for a parameter. \emph{(Made
       precise in Unit 7.)} \\
\end{tabularx}
\end{skillbox}

% fixedskillbox never splits, so the phase table stays intact on one page.
% THE PHASE TABLE IS A CONTRACT: 5 / 35 / 10 / 10 — the I do is ~20 of the 35 and the Guided
% Practice ~14; the debrief is spoken; the homework start is the second formative read.
\begin{fixedskillbox}[Lesson at a Glance --- \MeetingLength]{frost}
\renewcommand{\arraystretch}{1.25}
\begin{tabularx}{\linewidth}{@{}l c X X@{}}
\textbf{Phase} & \textbf{Min} & \textbf{Students} & \textbf{Teacher} \\
\hline
Warm-Up & 5 & Last night's market arithmetic, then the two groups hiding in the paragraph ---
  silent & Debrief item~3 aloud; write ``$50$ asked, $800$ reported on'' on the board and leave
  it up all period \\
Guided Notes \& Practice & 35 & Fill the vocabulary box as each term is named; work rows
  1--4 of the table (problems 1--15) with the pen in their hand; then the Guided Practice row,
  \emph{Riverbend Jobs} (16--19), together & \textbf{I do / we do, row by row}: read the
  definition, mark up the display, work problems 1, 5, 10, 12; the class works the rest (20)
  $\to$ \textbf{we do} Guided Practice 16--19 (15) \\
Debrief & 10 & Pens down, papers up --- answer aloud, nothing to fill in & Put Ana, Ben, and
  Cleo back on the board and take the ``somebody counted wrong'' reflex, then the
  ``a percent is a statistic'' reflex \\
Close \& Assign & 10 & \textbf{Start the homework in class}, alone and silent & Announce the
  homework and its form (packet or Desmos), circulate for the second formative read, preview
  Lesson~1.3 \\
\end{tabularx}
\end{fixedskillbox}

\begin{fixedskillbox}[Warm-Up --- Activate Prior Knowledge (5 min)]{frost}
\begin{tabularx}{\linewidth}{@{}X X@{}}
\begin{minipage}[t]{\linewidth}\vspace{0pt}
    \textbf{The three items and what each seeds}
    \begin{itemize}[leftmargin=1.1em, itemsep=1pt, topsep=2pt]
      \item \textbf{1.} $26$ of $50$ as a percent. \textbf{Spirals} to Lesson 1.1 --- and it is
            the $52\%$ that problem~5 will label a \emph{statistic}.
      \item \textbf{2.} $52\%$ of $800$. \textbf{Seeds:} the estimate --- row~2 (\emph{Parameter \&
            Statistic}) says you use the statistic you have to estimate the parameter you want,
            and problem~19 calls the $360$ an estimate, not a count.
      \item \textbf{3.} Two groups, no vocabulary. \textbf{Seeds:} population and sample ---
            row~1 attaches the names to a distinction students have already drawn.
    \end{itemize}
\end{minipage}
&
\begin{minipage}[t]{\linewidth}\vspace{0pt}
    \textbf{Running it}
    \begin{itemize}[leftmargin=1.1em, itemsep=1pt, topsep=2pt]
      \item \textbf{Debrief item~3 aloud} and write the two counts on the board --- \emph{$50$
            asked, $800$ reported on} --- in a student's own words. Leave it up all period;
            row~1 hangs \emph{sample} and \emph{population} on that line.
      \item Item~3's last question is the tell. \emph{The $52\%$} was counted; the $416$ was
            computed from it. A student who says $416$ has confused a prediction with a count
            and needs row~1 slowly.
      \item Items 1--2 are last night's arithmetic --- say so, so nobody thinks they are being
            retested. Do not let them expand.
    \end{itemize}
\end{minipage}
\end{tabularx}
\end{fixedskillbox}

\boxguard
\begin{skillbox}[Hook --- before anyone sits down]{frost}
The hook is on \textbf{page 1 of the notes}, under the vocabulary box, and on the board: \textit{``Three volunteers. The same Saturday morning. Each asked $50$ shoppers the same question. Ana got $52\%$. Ben got $48\%$. Cleo got $58\%$.''} \textbf{Ask: who made the mistake?} Students circle a name (or \emph{nobody}) and write one line of reason in the hook box \emph{before} anyone speaks; then take the hand vote and write the split on the board. Most rooms vote for somebody, or split the difference (``the real answer is $52.7\%$'').
\textbf{Do not settle it.} Say only that the answer is problem~14, in the \emph{One
Parameter} row, and that they will vote again there. The vote is what makes the crux land: students have to have
committed to \emph{somebody is wrong} before the notes take it away.
\end{skillbox}

\boxguard
\begin{skillbox}[Guided Notes \& Practice --- I do, then we do (35 min)]{frost}
\textbf{This is the whole lesson.} There is no group activity, and there is no independent
practice set on this page: \textbf{the homework is the individual practice}, started in class.
The notes are \textbf{one Main Ideas / Notes table}: four instruction rows (problems 1--15),
then the Guided Practice row (16--19). \textbf{The rhythm in every row is the same}: read the
definition aloud, mark up the display, work the first problem of the grid yourself, then hand
the rest to the room with the pen in their hand. The vocabulary box on page~1 is filled
\emph{as each term is named}, never in advance. The only blanks are the five cells of the
constraint table (problem~9); every other answer is written in a problem's own space.
\begin{multicols}{2}
\textbf{I do / we do, row by row (about 20 min).}

\emph{Row 1, Population \& Sample (problems 1--4).} Open on the line already on the board: $50$
asked, $800$ reported on. Read the definition, then have students name the parts of the
picture --- the box is the \emph{population}, the oval is the \emph{sample}, $n = 50$. Work
problem~1 yourself; the class does 2--4. Problem~3 is there because a population is not
always people; problem~4 is the census. Four minutes.

\emph{Row 2, Parameter \& Statistic (problems 5--8).} Read the two definitions and the memory
hook, then name the two numbers in the display --- the $52\%$ is a \emph{statistic}, the true
percent a \emph{parameter}. Work problem~5 yourself. \textbf{Take problem~8 last, and take a
hand count before anyone writes}: $61\%$ from district records. Half the room will say
\emph{statistic} because it is a percent. Then ask the only question that matters --- \emph{who
does it describe?} --- and let the room write the answer. Six minutes.

\emph{Row 3, Constraints (problems 9--11).} Give the five meanings from the cards, then take
the five situations of problem~9 from the class; the last row (\emph{access}) is the one they
will not have a word for. Work problem~10 yourself; 11 goes to the room. Four minutes.

\emph{Row 4, One Parameter (problems 12--15) --- the crux.} Put Ana, Ben, and Cleo up on the
display, work problem~12 live (Ben's and Cleo's statistics), have the class do 13 --- and then
\textbf{re-take the hook vote at problem~14 before you say anything}. Three statistics, one
parameter, nobody wrong. Land the printed sentence word by word: \emph{the statistic moves;
the parameter stays put}. Problem~15 (a fourth volunteer asks $200$) is the counterweight ---
closer, usually, but never exactly. Six minutes.

\columnbreak
\textbf{We do (about 15 min).} The Guided Practice row, \emph{Riverbend Jobs} (problems
16--19) --- the same moves the homework will ask alone, on a school instead of a market.
\textbf{Students hold the pen; you hold the questions.} Problem~16 goes on them cold (the
groups, $n$, and a constraint). Problem~17 is warm-up arithmetic; the labels are the point ---
\emph{statistic}, then \emph{estimate of a parameter}. \textbf{Problem~18 is the crux
transferred}: a second $50$ students, $46\%$ instead of $40\%$. Take a hand count on \emph{did
somebody make a mistake} before anyone writes, and make both sides defend it. Problem~19 is the
over-claim: \emph{``Exactly $360$\ldots''} --- not allowed, and the rewrite is the sentence the
homework will ask for.

Circulate during Guided Practice and demand the reason, not the word:
\begin{itemize}[leftmargin=1.1em, nosep]
  \item \textbf{Problem 16:} ``Who is the question about? Now --- who did she actually ask?''
  \item \textbf{Problem 17:} ``You wrote \emph{statistic} for the $360$. Who did she count to
        get $360$?'' (Nobody --- it was computed from the $40\%$.)
  \item \textbf{Problem 18:} ``Which of the two surveys found the true percent?'' A student who
        picks one has not left the hook vote.
  \item \textbf{Problem 19:} ``Who counted $360$ students?'' Nobody --- it was computed from
        the $40\%$. An estimate is not a count.
\end{itemize}
While circulating, pick the papers to put up in the debrief --- one problem~18 that said
\emph{neither, one parameter}, and one that picked a survey. \textbf{Your formative read comes
twice}: here, and again in the last ten minutes as students start the homework.
\end{multicols}
\end{skillbox}

\boxguard[24]
\begin{skillbox}[Debrief --- whole class, spoken (10 min)]{frost}
Pens down, papers up. \textbf{This debrief is spoken} --- there is no box at the end of the
notes to fill in, so it lives entirely on the board and in the room.
\begin{multicols}{2}
\begin{enumerate}[leftmargin=1.3em, nosep, itemsep=3pt]
  \item \textbf{Put Ana, Ben, and Cleo back up} and make a student say the sentence: three
        statistics, one parameter, nobody wrong. Then state it as PS.DC.1e: the value of a
        statistic varies from sample to sample. That is sampling variability, and it is the
        reason Unit~7 exists.
  \item Put up the two problem~18s you picked. The difference between them is
        not arithmetic --- it is whether the student still believes one survey ``won.'' Read
        both aloud and let the room hear which one a board member could act on.
  \item Put problem~8 back up and take the \emph{a-percent-is-a-statistic} reflex:
        \emph{``Who does the $61\%$ describe? Then what is it?''}
  \item Take the over-claim last, from problem~19: \emph{``Exactly $360$ Riverbend students have a part-time job.''} What word is wrong, and why (\emph{exactly} --- the $360$ is an estimate, and the true count is the parameter nobody knows).
\end{enumerate}
\columnbreak
\textbf{Check for understanding --- ask it cold, whole class.} \emph{``A school of $1{,}200$
students. The yearbook staff asks $60$ students and finds $45\%$ have a part-time job; the
newspaper asks a different $60$ and finds $50\%$. A teacher says one of them must have made
an error. Is the teacher right? What is the true percent?''} Hands down, thirty seconds of
think time, then cold-call three students. Correct answer: \textbf{no} --- two samples, two
statistics, sampling variability; the true percent is the one parameter for all $1{,}200$,
and neither survey found it. \emph{This replaces the exit ticket.}

\textbf{Formative read.} Sort the three answers as they come: \textbf{(a)} said \emph{no} and
named one parameter; \textbf{(b)} said \emph{no} but guessed ``the real answer is halfway'';
\textbf{(c)} picked a survey or said someone erred. Pile (b) is the teachable one and will be
large --- name it as an interesting guess Unit~7 makes precise, not as nonsense. \textbf{If
pile (c) runs past a third of the room, open Lesson~1.3 by re-running the \emph{One Parameter} row
(problems 12--15)}, and say so plainly.

\textbf{If the debrief runs short}, cut items 3 and~4 --- item~1 is the one that cannot be
cut. \textbf{Do not let the debrief eat the last ten minutes} --- the homework start is not
optional padding, it is your second formative read.
\end{multicols}
\end{skillbox}

\boxguard
\begin{skillbox}[Homework --- scored, started in class, due after two study halls]{goldbox}
Six exercises on the \textbf{Harbor Point Community Pool} ($1{,}500$ members; $75$ asked on a
Tuesday evening): \textbf{1} name the population, sample, and sample size; \textbf{2} the
frequency table with relative frequencies (\emph{the spiral to Lesson 1.1}); \textbf{3} scale
$36\%$ to $1{,}500$ and write the sentence \emph{as an estimate}; \textbf{4} label four numbers
--- rows (ii) and (iv) are the \textbf{contrast pair}: two percents, one a statistic and one a
parameter drawn from front-desk records; \textbf{5} name four constraints; \textbf{6} a second
sample of $75$ gives $40\%$ against the first $36\%$ --- \emph{the crux head-on}. The board's
over-claiming sentence lives in problem~19 of the notes, not here: the homework is capped at
\textbf{two pages} (user decision 2026-09-06).

\textbf{Start it in the last ten minutes of the period, in class and alone.} \textbf{Scoring:}
12 points --- 2 per item. \textbf{Item~6(b) is where the misconception shows}: a student who
explains the difference by ``the second staff member was more careful'' has learned the word
and not the rule. \textbf{Item~4 row (iv) carries the second check.}

\textbf{Packet or Desmos --- decided here.} The packet pages are authored and are the default.
Desmos Classroom has percent-of-a-total practice that would cover items 2--3, but nothing that
asks a student to explain why two surveys disagree or what a sentence over-claims --- the items
this lesson exists for --- so \textbf{1.2 is a packet night}. Say so aloud at Close \& Assign.

\textbf{Preview of Lesson~1.3.} The staff member asked whoever showed up on a Tuesday evening
--- lap swimmers keep evening hours. Do not resolve that here; it is 1.3's opening (how a sample
gets chosen) and 1.4's bias work.

\textbf{Connections \& Big Ideas.} \textit{Unit 1 core idea:} a conclusion is only as
trustworthy as the study that produced it. \textit{Standards covered:} PS.DC.1d, PS.DC.1e
(with PS.DC.1b carried forward in the conclusion items). \textit{Spirals forward to:} sampling
techniques (1.3, PS.DC.2a--b), bias (1.4, PS.DC.2c), the full-cycle project (1.8), and ---
directly --- sampling distributions and margin of error (7.1--7.2, PS.IS.1a, d).
\textit{Reaches back to:} Lesson 1.1's data cycle and conclusion in context, and Algebra~1
percent and proportional reasoning.
\end{skillbox}

\boxguard
\begin{skillbox}[Watch For (while circulating)]{redbox}
\begin{itemize}[leftmargin=1.2em, nosep]
  \item \textbf{``Somebody counted wrong''} (problem~14, problem~18, HW~6(b) --- the target
        misconception). Probe: \emph{``How many true percents are there for the whole
        population? Then how many of the three found it?''}
  \item \textbf{``The real answer is in the middle''} (problem~15, the cold check). Not wrong,
        not yet right: name it as the guess Unit~7 makes precise. Probe: \emph{``Would a fourth
        volunteer land on your middle number?''}
  \item \textbf{A percent is a statistic; a count is a parameter} (problem~8, HW~4(iv)).
        Probe: \emph{``Who does the number describe --- the ones asked, or everyone?''}
  \item \textbf{The estimate reported as a count} (problems~17 and~19, HW~3): ``$360$
        students have a job.'' Probe: \emph{``Who did she count to get $360$?''}
  \item \textbf{Population named as whoever was handy} (problems~1--3 and~16, HW~1): ``the
        population is the $50$ students.'' Probe: \emph{``Who is the question \emph{about}?''}
  \item \textbf{Over-claiming beyond the sample} (problem~19, HW~6(c)): from $75$ on a Tuesday to
        every member, and from \emph{what do you use the pool for} to \emph{do you want a
        second lane}. Probe: \emph{``Who was asked, and what were they asked?''}
\end{itemize}
Cold-call prompts: \emph{``Name a parameter about this school that nobody has ever
computed.''} \quad \emph{``Give me a study where the population is not people.''}
\end{skillbox}

\boxguard
\begin{skillbox}[Close \& Assign (10 min)]{goldbox}
\textbf{Students start the homework here, in class, alone and silent --- this is the individual
practice, and the first minutes of it are your best formative read.} Hand the launch first ---
six exercises on paper, scored out of 12, due at the start of the first class after two study halls --- and point out that the
\emph{Keep in Mind} box on the cover is the page to flip back to. Circulate: \textbf{item~6(b)
is the column that tells you whether the \emph{One Parameter} row landed}, and it is on page~2, so put students
on 1, 4, and~6 first and let 2, 3, and~5 go home. Sort what you see into three piles as you
walk --- said \emph{no mistake} and named one parameter (ready); said \emph{no mistake} with
no reason (ready, needs the language); blamed a staff member or picked a survey (the
misconception survived) --- and open Lesson~1.3 by re-running problems 12--15 if that
third pile ran past a third of the room. Then close on what changed: \emph{``You walked in
knowing how to turn $26$ of $50$ into a percent. You walk out knowing that percent is a
statistic, that the number you actually want is a parameter you will never see, and that two
honest samples disagreeing is not a mistake --- it is the whole reason the rest of this course
exists.''} \textbf{Preview:} Lesson~1.3, \emph{Choosing a Sample} --- today three volunteers
each grabbed $50$ shoppers; next time we ask \emph{which} $50$, and find out that how you
choose them decides whether the statistic is worth anything.
\end{skillbox}

% ── Teacher notes — one per component, in packet order (three) ──────────────
% Teacher-only prose lives HERE, never in a _key: a note is the one block in a key with no
% counterpart in the blank. No note for the debrief or the homework start — both are fully
% specified in their own boxes above.
\begin{teachernote}[Warm-Up]
Five minutes, silent, timed. Items 1 and~2 are last night's homework arithmetic on purpose ---
say so, so nobody thinks they are being retested. \textbf{Item~3 is the one to read while
collecting}: students who write $50$ and $800$ in the right rows and circle the $52\%$ have
already made today's distinction without the vocabulary, and row~1 can move fast.
Students who put $416$ in a row, or name it as the counted number, have confused an estimate
with a group and will need row~1 slowly. Write the two counts on the board in a student's
own words and \textbf{leave them up all period}; you will hang \emph{sample} and
\emph{population} on that line in row~1 and circle it again in the debrief. Do \emph{not}
supply the words \emph{sample} and \emph{population} here --- the warm-up is deliberately
wordless about them.
\end{teachernote}

\begin{teachernote}[Guided Notes \& Practice]
\textbf{This one document is the lesson --- pace it 20 / 15, then 10 for the debrief, then 10 for
the homework start.} It is one table, and every row runs the same way: you read the
definition, mark up the display, and work the first problem of the grid (1, 5, 10, 12); the
class works the rest with the pen in their hand. Rows 1 and~2 must move: four minutes on the
two groups, six on parameter and statistic. \textbf{Take the hand count on problem~8 before
anyone writes.} If students write first, the trap is gone; if they commit to \emph{statistic}
and \emph{then} ask who the number describes, the rule arrives as the answer to a question they
already have.

\textbf{Spend the time in row~4, \emph{One Parameter}}, which carries the misconception. Give
the five constraint cards and take problem~9 from the class quickly --- four minutes for row~3
--- then slow down. Work problem~12 live (Ben's and Cleo's statistics), let the class do~13,
and \textbf{re-take the hook vote at problem~14 before you say anything}: Ana, Ben, Cleo, or
nobody. Let the room argue. Then land the printed sentence word by word. Do not resolve it by
repeating the definition; count the parameters out loud --- \emph{one} --- and count the
statistics --- \emph{three}. Problem~15 is the counterweight: a bigger sample lands closer,
usually, and never exactly.

\textbf{Do not skip problem~19} --- the over-claim and its rewrite. Without it students leave
believing an estimate is a count, which is the version of this error that resurfaces in every
conclusion they write for the rest of the unit.

\textbf{Guided Practice (16--19) is where the pen changes hands.} \emph{Riverbend Jobs} runs
the four moves the homework will ask alone --- name the groups and the constraint, compute and
label, the second sample, the report sentence --- on a school instead of a market. Problem~18
is the crux transferred: $46\%$ against $40\%$, and a student who picks a survey has not left
the hook vote. If you only get through 16, 17, 18, and~19's verdict, that is the right trade;
the rewrite can be said aloud.

\textbf{There is no independent practice set on this page, and that is deliberate --- the
homework is the individual practice.} Guided Practice is the last row; the period ends with
students starting the homework in front of you, which is where your second formative read
comes from. Item~6(b) is the one to watch. If the ``somebody counted wrong'' pile is more than
a third of the room, \textbf{open Lesson~1.3 by re-running problems 12--15}, and say so
plainly.

\textbf{Debrief (10 min), spoken.} Ana, Ben, and Cleo back on the board with a student saying
\emph{three statistics, one parameter}; it is the one move that cannot be cut. \textbf{Do not
let the debrief eat the last ten minutes} --- the homework start is not optional padding.
\end{teachernote}

\begin{teachernote}[Homework]
\textbf{Scored, 12 points (2 per item), due at the start of the first class after two study halls.} The ten minutes at the end of the
period are a \emph{start}, not a completion: put students on 1, 4, and~6 while you can still
catch a wrong start, and let 2, 3, and~5 go home. Item~6 is the one to read first when
scoring --- any student who explains $36\%$ versus $40\%$ by someone being ``more careful'' has
learned the word and not the rule. Item~4 is the fullest diagnostic in the set: row (iv) is a
percent computed from the records of all $1{,}500$ members, and a student who labels it a
statistic is still sorting by what the number looks like. Score items 1--5 for correctness and
item~6 for the \emph{reason}, not the verdict. \textbf{When you score the page, sort item~6
into three piles} --- said \emph{no} and named one parameter (ready); said \emph{no} vaguely
(ready, needs the language); blamed a staff member or picked a survey (the misconception
survived). If more than a third land in the last pile, reopen Ana, Ben, and Cleo in the
Lesson~1.3 warm-up debrief rather than letting it ride.

\textbf{Packet is the call for 1.2.} Desmos carries the percent arithmetic and nothing else in
this set; do not swap it in. Reopen the Tuesday-evening sample at the start of Lesson~1.3 --- it
is 1.3's opening.
\end{teachernote}
\end{document}
